% ============================================================
% Round 05: Diffusion-DPO / Image Editing Preference Optimization
% ============================================================

\subsection{Round 05：Diffusion-DPO 与图像编辑偏好优化}

\subsubsection{Highlight}

\begin{conceptbox}
\textbf{本章相关脉络：}
\begin{itemize}[leftmargin=1.5em]
  \item \textbf{DPO}：arXiv 2023-05-29，Direct Preference Optimization，把 RLHF 的 reward + PPO 过程改写成 pairwise classification loss。论文：\url{https://arxiv.org/abs/2305.18290}
  \item \textbf{Pick-a-Pic}：arXiv 2023-05-02，真实用户 text-to-image 偏好数据，包含 prompt 和成对图像偏好。论文：\url{https://arxiv.org/abs/2305.01569}，代码/数据：\url{https://github.com/yuvalkirstain/PickScore}
  \item \textbf{Diffusion-DPO}：arXiv 2023-11-21，把 DPO 改写到 diffusion model，用 diffusion ELBO / denoising error 近似图像 likelihood。论文：\url{https://arxiv.org/abs/2311.12908}，代码：\url{https://github.com/SalesforceAIResearch/DiffusionDPO}
  \item \textbf{Diffusion-SDPO}：arXiv 2025-11-05，后续改进，指出标准 Diffusion-DPO 可能在扩大偏好间隔时破坏 winner 分支，并提出 safeguarded loser update。论文：\url{https://arxiv.org/abs/2511.03317}，代码：\url{https://github.com/AIDC-AI/Diffusion-SDPO}
\end{itemize}
\end{conceptbox}

\subsubsection{本章要解决的问题}

上一章我们讲了 reward：

\[
R(I_\text{src},c,I_\text{edit})
\]

但 reward model + RL 不是唯一后训练路线。

如果我们有偏好数据：

\[
(I_\text{src},c,I_w,I_l)
\]

其中：

\begin{itemize}[leftmargin=1.5em]
  \item $I_w$ 是 chosen / winner，也就是更好的编辑结果；
  \item $I_l$ 是 rejected / loser，也就是较差的编辑结果；
  \item 二者对应同一个原图 $I_\text{src}$ 和同一个编辑指令 $c$。
\end{itemize}

那么可以直接训练模型：

\[
\text{提高 } I_w \text{ 的相对概率}
\quad
\text{降低 } I_l \text{ 的相对概率}
\]

这就是 DPO / Diffusion-DPO 的核心。

\begin{intubox}
\textbf{直觉：} SFT 只告诉模型“这个答案可以学”；DPO 额外告诉模型“在这两个答案里，这个更好，那个更差”。
\end{intubox}

\subsubsection{先复习文本 DPO}

文本 DPO 的数据是：

\[
(x,y_w,y_l)
\]

其中 $y_w$ 是 chosen answer，$y_l$ 是 rejected answer。

文本模型的 DPO loss 是：

\[
\mathcal{L}_\text{DPO}
=
-
\log
\sigma
\left(
\beta
\left[
\log\frac{\pi_\theta(y_w\mid x)}{\pi_\text{ref}(y_w\mid x)}
-
\log\frac{\pi_\theta(y_l\mid x)}{\pi_\text{ref}(y_l\mid x)}
\right]
\right)
\]

把它展开：

\[
\mathcal{L}_\text{DPO}
=
-
\log
\sigma
\left(
\beta
\left[
\log \pi_\theta(y_w\mid x)
-
\log \pi_\text{ref}(y_w\mid x)
-
\log \pi_\theta(y_l\mid x)
+
\log \pi_\text{ref}(y_l\mid x)
\right]
\right)
\]

这个式子的含义是：

\[
\text{新模型相对 reference，要更偏向 chosen，而不是 rejected}
\]

其中 $\beta$ 控制偏好优化强度。$\beta$ 越大，模型越激进地拉开 chosen 和 rejected；$\beta$ 越小，模型越保守。

\begin{warnbox}
\textbf{reference 的作用：} DPO 不是无约束地提高 chosen 概率。它比较的是 $\pi_\theta$ 相对 $\pi_\text{ref}$ 的变化，这等价于在偏好优化里保留一个 KL regularization 的锚点，防止模型漂太远。
\end{warnbox}

\subsubsection{迁移到图像：直接难点是什么}

图像编辑偏好数据可以写成：

\[
x=(I_\text{src},c)
\]

\[
(x,I_w,I_l)
\]

如果完全照搬文本 DPO，应该写成：

\[
\mathcal{L}_\text{edit-DPO}
=
-
\log
\sigma
\left(
\beta
\left[
\log\frac{\pi_\theta(I_w\mid I_\text{src},c)}
{\pi_\text{ref}(I_w\mid I_\text{src},c)}
-
\log\frac{\pi_\theta(I_l\mid I_\text{src},c)}
{\pi_\text{ref}(I_l\mid I_\text{src},c)}
\right]
\right)
\]

问题是：diffusion model 通常不像 LLM 那样直接给：

\[
\log \pi_\theta(I\mid I_\text{src},c)
\]

LLM 可以把整段文本 logprob 写成 token logprob 之和：

\[
\log \pi_\theta(y\mid x)
=
\sum_i
\log \pi_\theta(y_i\mid x,y_{<i})
\]

但 diffusion 的最终图像概率需要把整条去噪轨迹积分 / 边缘化掉：

\[
p_\theta(z_0\mid c)
=
\int
p(z_T)
\prod_{t=1}^{T}
p_\theta(z_{t-1}\mid z_t,c)
\,
dz_{1:T}
\]

这个积分不好直接算。

\begin{warnbox}
\textbf{Diffusion-DPO 的核心难点：} 文本 DPO 需要 $\log \pi_\theta(y\mid x)$；diffusion 模型没有一个像 token logprob 那样直接好算的最终图像 logprob。
\end{warnbox}

\subsubsection{关键替换：用去噪误差近似 Logprob}

Diffusion-DPO 的做法是：不要直接算最终图像的精确 likelihood，而是用 diffusion ELBO / denoising objective 给出一个可训练的 surrogate。

先看普通 diffusion SFT 的噪声预测 loss。给定目标图像 latent：

\[
z_0
=
E_\text{VAE}(I)
\]

加噪得到：

\[
z_t
=
\sqrt{\bar{\alpha}_t}z_0
+
\sqrt{1-\bar{\alpha}_t}\epsilon,
\qquad
\epsilon\sim\mathcal{N}(0,I)
\]

模型预测噪声：

\[
\epsilon_\theta(z_t,t,x)
\]

去噪误差是：

\[
\ell_\theta(I,t,x)
=
\left\|
\epsilon
-
\epsilon_\theta(z_t,t,x)
\right\|^2
\]

直觉上：

\[
\ell_\theta \text{ 越小}
\quad\Rightarrow\quad
\text{模型越会生成这张图}
\]

所以可以粗略理解为：

\[
\log \pi_\theta(I\mid x)
\approx
\text{constant}
-
\sum_t
\omega_t
\ell_\theta(I,t,x)
\]

其中 $\omega_t$ 是不同时间步的权重。

\begin{intubox}
\textbf{直觉：} 对 diffusion 来说，“这张图概率高”可以近似理解成：“把这张图加噪到任意时间步后，模型都很擅长把它去噪回来”。
\end{intubox}

\subsubsection{Diffusion-DPO 的核心公式}

现在有一组图像偏好样本：

\[
x=(I_\text{src},c)
\]

\[
I_w \succ I_l
\]

分别编码成 latent：

\[
z_0^w=E_\text{VAE}(I_w),
\qquad
z_0^l=E_\text{VAE}(I_l)
\]

采样同一个时间步 $t$，并分别加噪：

\[
z_t^w
=
\sqrt{\bar{\alpha}_t}z_0^w
+
\sqrt{1-\bar{\alpha}_t}\epsilon^w
\]

\[
z_t^l
=
\sqrt{\bar{\alpha}_t}z_0^l
+
\sqrt{1-\bar{\alpha}_t}\epsilon^l
\]

对当前模型计算去噪误差：

\[
\ell_\theta^w
=
\left\|
\epsilon^w
-
\epsilon_\theta(z_t^w,t,x)
\right\|^2
\]

\[
\ell_\theta^l
=
\left\|
\epsilon^l
-
\epsilon_\theta(z_t^l,t,x)
\right\|^2
\]

对 reference model 也计算同样的误差：

\[
\ell_\text{ref}^w
=
\left\|
\epsilon^w
-
\epsilon_\text{ref}(z_t^w,t,x)
\right\|^2
\]

\[
\ell_\text{ref}^l
=
\left\|
\epsilon^l
-
\epsilon_\text{ref}(z_t^l,t,x)
\right\|^2
\]

因为 logprob 近似等于负的去噪误差，所以图像 DPO 的 margin 可以写成：

\[
m_\theta
=
\beta\omega_t
\left[
\left(
\ell_\theta^l-\ell_\text{ref}^l
\right)
-
\left(
\ell_\theta^w-\ell_\text{ref}^w
\right)
\right]
\]

于是 Diffusion-DPO loss 是：

\[
\mathcal{L}_\text{Diff-DPO}
=
-
\log
\sigma(m_\theta)
\]

把它读成中文就是：

\[
\text{chosen 的去噪误差，相对 ref 要变小}
\]

\[
\text{rejected 的去噪误差，相对 ref 可以变大}
\]

所以：

\[
\left(
\ell_\theta^w-\ell_\text{ref}^w
\right)
\downarrow
\]

\[
\left(
\ell_\theta^l-\ell_\text{ref}^l
\right)
\uparrow
\]

margin 就会变大，loss 就会变小。

\begin{summarybox}
\textbf{Diffusion-DPO 一句话：}\\
文本 DPO 比较的是 chosen / rejected 的 logprob ratio；Diffusion-DPO 比较的是 chosen / rejected 的 denoising error 相对 reference 的变化。因为 diffusion 的 logprob 可以用 ELBO / 去噪误差近似，所以 DPO loss 可以落到噪声预测 MSE 上。
\end{summarybox}

\subsubsection{为什么公式里是 Rejected 减 Chosen}

这个符号很容易看晕。我们从 logprob 角度推一次。

文本 DPO 的核心 margin 是：

\[
\beta
\left[
\log\frac{\pi_\theta(I_w\mid x)}{\pi_\text{ref}(I_w\mid x)}
-
\log\frac{\pi_\theta(I_l\mid x)}{\pi_\text{ref}(I_l\mid x)}
\right]
\]

现在用：

\[
\log\pi_\theta(I\mid x)
\approx
-
\omega_t\ell_\theta(I,t,x)
\]

代入 chosen：

\[
\log\frac{\pi_\theta(I_w\mid x)}{\pi_\text{ref}(I_w\mid x)}
\approx
-
\omega_t
\left(
\ell_\theta^w-\ell_\text{ref}^w
\right)
\]

代入 rejected：

\[
\log\frac{\pi_\theta(I_l\mid x)}{\pi_\text{ref}(I_l\mid x)}
\approx
-
\omega_t
\left(
\ell_\theta^l-\ell_\text{ref}^l
\right)
\]

两者相减：

\[
-
\omega_t
\left(
\ell_\theta^w-\ell_\text{ref}^w
\right)
-
\left[
-
\omega_t
\left(
\ell_\theta^l-\ell_\text{ref}^l
\right)
\right]
\]

整理后就是：

\[
\omega_t
\left[
\left(
\ell_\theta^l-\ell_\text{ref}^l
\right)
-
\left(
\ell_\theta^w-\ell_\text{ref}^w
\right)
\right]
\]

所以你看到的是：

\[
\text{rejected error improvement}
-
\text{chosen error improvement}
\]

更准确说，是：

\[
\text{rejected 相对 ref 的误差变化}
-
\text{chosen 相对 ref 的误差变化}
\]

由于 logprob 是负误差，所以符号会反过来。

\begin{intubox}
\textbf{记忆法：} 文本里要让 chosen logprob 更高；diffusion 里等价于让 chosen denoising error 更低。logprob 越高越好，error 越低越好，所以从 logprob 换成 error 时符号会翻转。
\end{intubox}

\subsubsection{图像编辑版 DPO 与 Text-to-Image DPO 的差别}

Text-to-image Diffusion-DPO 的条件通常是 prompt：

\[
x=c_\text{text}
\]

图像编辑 DPO 的条件是：

\[
x=(I_\text{src},c)
\]

因此每次计算 chosen / rejected 的去噪误差时，模型都必须看到同一个原图和同一个编辑指令：

\[
\epsilon_\theta(z_t^w,t,I_\text{src},c)
\]

\[
\epsilon_\theta(z_t^l,t,I_\text{src},c)
\]

这很关键。因为 chosen 好，不只是因为它漂亮，而是因为它相对于同一张原图做了更正确的编辑。

\begin{warnbox}
\textbf{图像编辑 DPO 的关键：} chosen / rejected 必须在同一个 source image 和同一个 instruction 下比较。否则模型学到的可能只是“哪张图更好看”，而不是“哪个编辑更正确”。
\end{warnbox}

\subsubsection{Preference Pair 可以从哪里来}

图像编辑偏好对可以来自几种方式。

\textbf{第一种：人工标注。}

给同一个编辑任务生成多个结果，让人类选择更好的：

\[
I_w \succ I_l
\]

人工偏好质量高，但成本高。

\textbf{第二种：reward / scorer 自动排序。}

用上一章的 reward：

\[
R(I_\text{src},c,I)
\]

对多个候选排序：

\[
R(I_a)>R(I_b)
\quad\Rightarrow\quad
I_a\succ I_b
\]

成本低，但会继承 reward model 的偏差。

\textbf{第三种：VLM judge。}

让 VLM 比较：

\[
(I_\text{src},c,I_a,I_b)
\]

输出哪个编辑更好。

这种方式更像自动化人类偏好标注，但也会有 judge bias。

\begin{summarybox}
\textbf{偏好数据的本质：}\\
Diffusion-DPO 不要求你有绝对 reward 分数，只要求你知道：在同一个编辑任务下，两个候选结果哪个更好。
\end{summarybox}

\subsubsection{DPO、SFT、Reward RL 的区别}

\begin{center}
{\small
\begin{tabular}{p{2.7cm}p{3.6cm}p{4.2cm}p{3.4cm}}
\hline
\textbf{方法} & \textbf{数据} & \textbf{训练信号} & \textbf{核心风险} \\
\hline
SFT
&
$(I_\text{src},c,I_\text{tgt})$
&
模仿 target image 的去噪 / 速度场
&
target 不一定代表人类最优偏好
\\
DPO
&
$(I_\text{src},c,I_w,I_l)$
&
提高 chosen 相对 rejected 的概率
&
受离线偏好数据覆盖范围限制
\\
Reward RL
&
采样结果 + reward
&
直接最大化 reward
&
训练不稳定、reward hacking
\\
\hline
\end{tabular}
}
\end{center}

\subsubsection{Diffusion-DPO 为什么不是普通 SFT}

如果只用 chosen 图像做 SFT：

\[
\mathcal{L}_\text{SFT}
=
\ell_\theta^w
\]

模型只知道：

\[
I_w \text{ 是可以学习的}
\]

但它不知道：

\[
I_w \text{ 比 } I_l \text{ 好在哪里}
\]

DPO 的训练信号是相对的：

\[
\mathcal{L}_\text{DPO}
=
-
\log
\sigma
\left(
\text{chosen score}
-
\text{rejected score}
\right)
\]

它要求模型拉开两者：

\[
\text{score}_\theta(I_w,x)
>
\text{score}_\theta(I_l,x)
\]

在 diffusion 里，score 来自负去噪误差：

\[
\text{score}_\theta(I,x)
\approx
-
\sum_t
\omega_t
\left(
\ell_\theta(I,t,x)-\ell_\text{ref}(I,t,x)
\right)
\]

\begin{intubox}
\textbf{直觉：} SFT 是“向好答案靠近”；DPO 是“相对于坏答案，更偏向好答案”。这就是为什么 DPO 需要 winner 和 loser 成对出现。
\end{intubox}

\subsubsection{Diffusion-DPO 的训练流程}

一个最小训练流程是：

\begin{enumerate}[leftmargin=1.5em]
  \item 从偏好数据中取 $(I_\text{src},c,I_w,I_l)$；
  \item 把 $I_w$ 和 $I_l$ 编码成 latents：$z_0^w,z_0^l$；
  \item 采样时间步 $t$；
  \item 分别对 winner / loser 加噪得到 $z_t^w,z_t^l$；
  \item 用当前模型计算 $\ell_\theta^w,\ell_\theta^l$；
  \item 用冻结 reference model 计算 $\ell_\text{ref}^w,\ell_\text{ref}^l$；
  \item 构造 DPO margin $m_\theta$；
  \item 最小化 $-\log\sigma(m_\theta)$。
\end{enumerate}

压缩成一行：

\[
(I_w,I_l)
\Rightarrow
(\ell_\theta^w,\ell_\theta^l)
\Rightarrow
\text{preference margin}
\Rightarrow
\text{DPO loss}
\]

\subsubsection{DPO 的优点和局限}

优点：

\begin{itemize}[leftmargin=1.5em]
  \item 不需要显式训练 reward model；
  \item 不需要在线 PPO / DDPO 那样每步采样再 policy gradient；
  \item 训练形式接近 supervised fine-tuning，工程上更稳定；
  \item 可以直接利用人工或自动构造的 preference pairs。
\end{itemize}

局限：

\begin{itemize}[leftmargin=1.5em]
  \item 它依赖离线偏好数据，数据覆盖不到的区域很难改；
  \item preference pair 的质量会直接决定训练方向；
  \item 如果 winner / loser 只反映美学差异，模型可能学不到编辑正确性；
  \item 过强的 DPO 可能破坏原模型的生成能力或 winner 分支质量；
  \item 它没有 RL 的在线探索能力。
\end{itemize}

\begin{warnbox}
\textbf{为什么有后续 SDPO / RL 改进：}\\
标准 Diffusion-DPO 的目标是拉大 chosen 和 rejected 的相对间隔，但“拉大间隔”不必然等于“chosen 变得更好”。如果更新主要通过让 loser 更差来实现，或者 winner 分支也被误伤，就可能出现训练退化。因此后续方法会加入更细的更新约束、dense reward 或在线 RL。
\end{warnbox}

\subsubsection{本章最小结论}

\begin{summarybox}
\textbf{Diffusion-DPO 的核心：}\\
文本 DPO 用 chosen / rejected 的 sequence logprob 做偏好分类；Diffusion-DPO 没有直接好算的 image logprob，于是用 diffusion ELBO / denoising error 作为 surrogate。图像编辑版 DPO 进一步把条件从 text prompt 扩展为 $(I_\text{src},c)$，让模型在同一个原图和指令下，更偏向 chosen 编辑结果，而不是 rejected 编辑结果。
\end{summarybox}

\subsubsection{本轮检查题}

\begin{enumerate}[leftmargin=1.5em]
  \item 文本 DPO 的 chosen / rejected logprob ratio 是什么？
  \item 为什么 diffusion model 不能像 LLM 一样直接计算最终图像的 token logprob？
  \item Diffusion-DPO 为什么可以用 denoising error 近似 logprob？
  \item 公式里为什么会出现 $\ell_\theta^l-\ell_\theta^w$ 这种 rejected 减 chosen 的结构？
  \item 图像编辑 DPO 和 text-to-image DPO 的条件输入有什么区别？
  \item DPO 相比 SFT 多用了什么信息？
  \item DPO 相比 reward RL 少了什么能力？
\end{enumerate}
